\documentclass[a4paper,11pt]{article}

\usepackage{amsmath,amssymb,amsthm,mathtools}
\usepackage{bm}
\usepackage{geometry}
\usepackage{enumitem}
\usepackage{hyperref}
\usepackage{array}

\geometry{
    top=25mm,
    bottom=25mm,
    left=25mm,
    right=25mm
}

%============================================================
% Theorem environments
%============================================================

\newtheorem{definition}{定義}[section]
\newtheorem{theorem}[definition]{定理}
\newtheorem{proposition}[definition]{命題}
\newtheorem{lemma}[definition]{補題}
\newtheorem{corollary}[definition]{系}
\newtheorem{remark}[definition]{注記}

%============================================================
% Notation
%============================================================

\newcommand{\Solver}{\mathcal{S}}
\newcommand{\SubSolver}{\preceq}
\newcommand{\Type}{\operatorname{Type}}
\newcommand{\DType}{\mathrm{D}}
\newcommand{\CType}{\mathrm{C}}
\newcommand{\MType}{\mathrm{M}}

\newcommand{\Aset}{\mathcal{A}}
\newcommand{\Rset}{\mathcal{R}}
\newcommand{\Th}{\mathrm{Th}}
\newcommand{\N}{\mathbb{N}}
\newcommand{\R}{\mathbb{R}}
\newcommand{\Pset}{\mathcal{P}}

\newcommand{\Eval}{\mathcal{G}}
\newcommand{\EvalMap}{\Psi}

%============================================================
% Title
%============================================================

\title{ECA Solver 集合に関する一般解の厳密解析 改訂版 v2 に基づく (A)--(H) の厳密証明}

\author{}
\date{}

\begin{document}

\maketitle
\tableofcontents

%################################################################
% 序論
%################################################################

\section{序論}

本稿では、拡張連続性公理
(Extended Continuity Axiom; ECA)
において構成された Solver 集合について、
改訂版 v2 の定義体系を基礎として、
(A)--(H) において導入された命題のうち、
実際に既存の数学および前節までの定義から導出可能なものを
厳密に証明する。

本稿では特に、

\begin{equation}
    \text{定義}
    \quad\Longrightarrow\quad
    \text{仮定}
    \quad\Longrightarrow\quad
    \text{導出可能な命題}
\end{equation}

という論理的階層を明確にする。

したがって、ECA における数学的対象として重要であっても、
現在の定義体系からは証明できない事項については、
それを無理に定理として扱わず、
必要な仮定または論考として明示する。

Solver を

\begin{equation}
    \Solver
    =
    \left\{
        S
        \,\middle|\,
        S=(\Aset_S,\Rset_S)
    \right\}
\end{equation}

とする。

ここで $\Aset_S$ は Solver $S$ における公理選択集合、
$\Rset_S$ は推論規則および構成制約の集合である。

さらに、Solver によって表現される理論を

\begin{equation}
    \mathcal{T}:
    \Solver\longrightarrow\Th
\end{equation}

によって表す。

一般評価空間を $X$ とし、
一般解を

\begin{equation}
    \Eval:
    \Solver\longrightarrow X
\end{equation}

とする。

評価空間から最終的な実数値評価を取り出す写像を

\begin{equation}
    \EvalMap:
    X\longrightarrow\R
\end{equation}

とすれば、

\begin{equation}
    g
    :=
    \EvalMap\circ\Eval
\end{equation}

によって Solver 集合上の最終評価関数を得る。

%################################################################
% (A)
%################################################################

\section{(A) Solver 集合の数学的構造}

\subsection{A.1 Solver の基本構造}

各 Solver $S\in\Solver$ を

\begin{equation}
    S=(\Aset_S,\Rset_S)
\end{equation}

とする。

ここで $\Aset_S$ は公理選択集合、
$\Rset_S$ は推論規則および構成制約である。

この定義から、Solver は単なる数値集合ではなく、
少なくとも公理選択および推論・構成規則からなる構造として
扱われる。

また、

\begin{equation}
    \mathcal{T}(S)
    =
    \mathcal{T}(\Aset_S,\Rset_S)
\end{equation}

とする。

\begin{remark}
$\mathcal{T}(S)$ は公理集合そのものではなく、
公理選択および推論規則・構成制約によって表現される理論である。
したがって、

\[
    \mathcal{T}(S)\neq\Aset_S
\]

である。
\end{remark}

%------------------------------------------------------------

\subsection{A.2 部分 Solver}

\begin{definition}[部分 Solver]
$S',S\in\Solver$ に対して

\begin{equation}
    S'\SubSolver S
\end{equation}

とは、$S'$ が $S$ の一部の要素から構成され、
かつ Solver として必要な構造を保持する独立した部分構造であることをいう。
\end{definition}

特に、

\begin{equation}
    S'\SubSolver S
    \Longrightarrow
    S'\subseteq S
\end{equation}

を要求する。

したがって、包含関係

\begin{equation}
    S'\subseteq S
\end{equation}

だけでは一般に

\begin{equation}
    S'\SubSolver S
\end{equation}

を導くことはできない。

%------------------------------------------------------------

\subsection{A.3 離散的要素と連続的要素}

Solver $S$ の要素 $u\in S$ について、
所定の数学的構造によって離散性および連続性を判定する。

離散的要素の集合を

\begin{equation}
    D_S
    :=
    \left\{
        u\in S
        \mid
        u\text{ が所定の離散的構造を持つ}
    \right\}
\end{equation}

とする。

同様に、

\begin{equation}
    C_S
    :=
    \left\{
        u\in S
        \mid
        u\text{ が所定の連続的構造を持つ}
    \right\}
\end{equation}

とする。

ここで重要なのは、
「非可算」と「連続」を無条件に同一視しないことである。

例えば、

\begin{equation}
    |S|\leq\aleph_0
\end{equation}

または

\begin{equation}
    |S|>\aleph_0
\end{equation}

という性質は集合論によって判定できるが、
後者だけから連続性を導くことはできない。

したがって、Solver の型判定には必要に応じて

\begin{equation}
    \text{集合論},
    \quad
    \text{位相論},
    \quad
    \text{順序論},
    \quad
    \text{代数学},
    \quad
    \text{解析学}
\end{equation}

等の既存数学を用いる。

%------------------------------------------------------------

\subsection{A.4 中間的要素}

\begin{definition}[中間的要素]
$u\in S$ が中間的であるとは、

\begin{equation}
    u\in D_S\cap C_S
\end{equation}

が成立することをいう。
\end{definition}

したがって、

\begin{equation}
    D_S\cap C_S\neq\varnothing
\end{equation}

ならば、Solver $S$ は少なくとも一つの
離散的かつ連続的な要素を含む。

これは単に

\begin{equation}
    \text{離散部分}\cup\text{連続部分}
\end{equation}

が存在するという意味ではない。

必要なのは、

\begin{equation}
    \exists u\in S:
    \quad
    u\in D_S
    \quad\land\quad
    u\in C_S
\end{equation}

である。

%------------------------------------------------------------

\subsection{A.5 Solver の型集合}

\begin{definition}[Solver の型集合]
Solver $S$ の型集合を

\begin{equation}
    \Type(S)
    \subseteq
    \{\DType,\CType,\MType\}
\end{equation}

とする。

\end{definition}

この定義によって、

\begin{equation}
    \Type(S)=\{\DType,\MType\}
\end{equation}

や

\begin{equation}
    \Type(S)
    =
    \{\DType,\CType,\MType\}
\end{equation}

のような複合型を許容する。

したがって、Solver を必ず一つの型に分類する必要はない。

%------------------------------------------------------------

\subsection{A.6 部分 Solver による型判定}

部分 Solver の型まで含めた型集合を

\begin{equation}
    \Type_{\mathrm{sub}}(S)
    :=
    \left\{
        \tau
        \,\middle|\,
        \exists S'\SubSolver S,\;
        \tau\in\Type(S')
    \right\}
\end{equation}

と定義する。

したがって、

\begin{equation}
    \Type_{\mathrm{sub}}(S)
    =
    \{\DType,\CType,\MType\}
\end{equation}

ならば、$S$ の内部には三種類の Solver 構造が存在する。

\begin{remark}
ここで $\Type(S)$ と
$\Type_{\mathrm{sub}}(S)$ は区別される。

前者は Solver 自身の構造を表し、
後者はその内部に含まれる部分 Solver まで調べた結果を表す。
\end{remark}

%------------------------------------------------------------

\subsection{A.7 既存数学による型判定}

Solver の型判定に既存数学が利用できる場合には、
ECA 固有の量よりも既存数学を優先する。

例えば濃度については、

\begin{equation}
    |S|\leq\aleph_0
\end{equation}

または

\begin{equation}
    |S|>\aleph_0
\end{equation}

によって可算性・非可算性を判定できる。

また位相構造が与えられている場合には、

\begin{equation}
    (S,\tau_S)
\end{equation}

について既存の位相論による性質を利用できる。

したがって、

\begin{equation}
    \boxed{
    \text{既存数学による判定}
    \longrightarrow
    \text{必要な場合のみ ECA の量を導入}
    }
\end{equation}

という原則を採用する。

%------------------------------------------------------------

\subsection{A.8 非可算性の継承}

\begin{proposition}[部分 Solver による非可算性の継承]
$S',S\in\Solver$ に対して

\begin{equation}
    S'\SubSolver S
\end{equation}

かつ

\begin{equation}
    |S'|>\aleph_0
\end{equation}

ならば、

\begin{equation}
    |S|>\aleph_0
\end{equation}

である。
\end{proposition}

\begin{proof}
$S'\SubSolver S$ より、

\begin{equation}
    S'\subseteq S
\end{equation}

である。

したがって包含写像

\begin{equation}
    \iota:S'\hookrightarrow S
\end{equation}

は単射であり、

\begin{equation}
    |S'|\leq|S|
\end{equation}

を得る。

仮定より

\begin{equation}
    |S'|>\aleph_0
\end{equation}

であるから、

\begin{equation}
    |S|
    \geq
    |S'|
    >
    \aleph_0.
\end{equation}

よって

\begin{equation}
    |S|>\aleph_0
\end{equation}

である。
\end{proof}

\begin{remark}
この命題は「任意の Solver が非可算である」ことを示すものではない。

示しているのは、

\[
\text{非可算な部分 Solver の存在}
\Longrightarrow
\text{全体 Solver の非可算性}
\]

という継承関係である。
\end{remark}

%################################################################
% (B)
%################################################################

\section{(B) 一般評価空間の構成}

\subsection{B.1 評価空間と評価写像}

Solver の評価構造を保持する空間を $X$ とする。

評価写像を

\begin{equation}
    \Phi:
    \Solver\longrightarrow X
\end{equation}

とする。

また、ECA における一般解を

\begin{equation}
    \Eval:
    \Solver\longrightarrow X
\end{equation}

と定義する。

ここで $\Phi$ と $\Eval$ は型が同一であっても、
概念上同一であるとは限らない。

%------------------------------------------------------------

\subsection{B.2 評価構造の表現}

Solver $S$ から得られる評価情報を

\begin{equation}
    \mathcal E(S)
\end{equation}

とする。

評価空間 $X$ は少なくとも

\begin{equation}
    \Phi(S)\in X
\end{equation}

となるように構成され、

\begin{equation}
    \Phi(S)
    =
    \operatorname{EvalStruct}(S)
\end{equation}

と表現できる場合を考える。

\begin{remark}
この時点では、$\operatorname{EvalStruct}$ の存在および
一意性を ECA の定義のみから証明しているわけではない。

これは評価空間を構成するための要件である。
\end{remark}

%------------------------------------------------------------

\subsection{B.3 ECA の評価量}

ECA において

\begin{equation}
    B,\qquad I,\qquad K
\end{equation}

が評価量として与えられる場合には、

\begin{equation}
    X
    =
    X_B\times X_I\times X_K
\end{equation}

と構成できる場合がある。

この場合、

\begin{equation}
    \Phi(S)
    =
    \left(
        B(S),I_S(S),K(S)
    \right)
\end{equation}

と表現できる。

ここで $I_S(S)$ は Solver の評価量としての $I$ であり、
公理選択空間 $I$ とは区別する。

%------------------------------------------------------------

\subsection{B.4 離散評価構造}

ECA の既存定義

\begin{equation}
    B(i)\subseteq i\subseteq S
\end{equation}

を用いる。

対応する公理選択を $i_S$ とすれば、

\begin{equation}
    E_{\mathrm d}(S)
    =
    \sum_{b\in B(i_S)}w_S(b)
\end{equation}

という離散評価を構成できる。

ただし、

\begin{equation}
    i_S
\end{equation}

の一意な復元については (H) において扱う。

%------------------------------------------------------------

\subsection{B.5 連続評価構造}

測度 $\mu_S$ が与えられ、
公理選択 $i_S$ が可測集合である場合、

\begin{equation}
    E_{\mathrm c}(S)
    =
    \mu_S(i_S)
    =
    \int_S
    \mathbf 1_{i_S}(u)\,d\mu_S(u)
\end{equation}

と表現できる。

ここで、

\begin{equation}
    \mathbf1_{i_S}:S\to\{0,1\}
\end{equation}

は指示関数である。

%------------------------------------------------------------

\subsection{B.6 中間 Solver の評価空間}

中間 Solver については、

\begin{equation}
    D_S\cap C_S\neq\varnothing
\end{equation}

が成立する。

必要に応じて評価空間を

\begin{equation}
    X
    =
    X_{\mathrm d}
    \times
    X_{\mathrm c}
    \times
    X_{\mathrm{dc}}
\end{equation}

と構成することができる。

ただし、これは中間 Solver の一般定義そのものではなく、
双方の評価構造を同一の評価空間に保持するための
具体的構成である。

%------------------------------------------------------------

\subsection{B.7 最終評価}

評価空間から実数値を取り出す写像を

\begin{equation}
    \EvalMap:X\to\R
\end{equation}

とする。

例えば、

\begin{equation}
    X=X_B\times X_I\times X_K
\end{equation}

かつ

\begin{equation}
    x=(B,I,K)
\end{equation}

ならば、

\begin{equation}
    \EvalMap(B,I,K)=BIK
\end{equation}

とできる。

したがって、

\begin{equation}
    g
    =
    \EvalMap\circ\Eval
\end{equation}

であり、

\begin{equation}
    g(S)
    =
    \EvalMap(\Eval(S))
\end{equation}

となる。

%################################################################
% (C)
%################################################################

\section{(C) 評価関数の測度論的性質}

\subsection{C.1 可測構造}

評価空間 $X$ に $\sigma$-代数 $\Sigma_X$ を与え、

\begin{equation}
    (X,\Sigma_X)
\end{equation}

を可測空間とする。

同様に Solver 集合上に

\begin{equation}
    \Sigma_{\Solver}
\end{equation}

を与え、

\begin{equation}
    (\Solver,\Sigma_{\Solver})
\end{equation}

を可測空間とする。

ここで、

\begin{equation}
    \Sigma_X
    \neq
    \mu_X
\end{equation}

であり、

\begin{equation}
    \text{可測構造}
    \neq
    \text{測度}
\end{equation}

である。

%------------------------------------------------------------

\subsection{C.2 評価写像の可測性}

\begin{proposition}[一般解の可測性]
もし

\begin{equation}
    \Eval:
    (\Solver,\Sigma_{\Solver})
    \longrightarrow
    (X,\Sigma_X)
\end{equation}

および

\begin{equation}
    \EvalMap:
    (X,\Sigma_X)
    \longrightarrow
    (\R,\mathcal B(\R))
\end{equation}

が可測ならば、

\begin{equation}
    g=\EvalMap\circ\Eval
\end{equation}

は可測である。
\end{proposition}

\begin{proof}
任意の

\begin{equation}
    A\in\mathcal B(\R)
\end{equation}

を取る。

$\EvalMap$ の可測性より、

\begin{equation}
    \EvalMap^{-1}(A)\in\Sigma_X.
\end{equation}

さらに $\Eval$ の可測性より、

\begin{equation}
    \Eval^{-1}
    \left(
        \EvalMap^{-1}(A)
    \right)
    \in\Sigma_{\Solver}.
\end{equation}

一方、

\begin{align}
    g^{-1}(A)
    &=
    (\EvalMap\circ\Eval)^{-1}(A)
    \\
    &=
    \Eval^{-1}
    \left(
        \EvalMap^{-1}(A)
    \right).
\end{align}

したがって、

\begin{equation}
    g^{-1}(A)\in\Sigma_{\Solver}.
\end{equation}

$A$ は任意だったので $g$ は可測である。
\end{proof}

%------------------------------------------------------------

\subsection{C.3 非負性}

\begin{proposition}[非負性]
もし

\begin{equation}
    \EvalMap(X)\subseteq[0,\infty)
\end{equation}

ならば、

\begin{equation}
    g(S)\geq0
    \qquad
    (\forall S\in\Solver)
\end{equation}

である。
\end{proposition}

\begin{proof}
任意の $S\in\Solver$ に対して

\begin{equation}
    \Eval(S)\in X
\end{equation}

である。

仮定より、

\begin{equation}
    \EvalMap(\Eval(S))\geq0.
\end{equation}

ところが、

\begin{equation}
    g(S)=\EvalMap(\Eval(S))
\end{equation}

であるから、

\begin{equation}
    g(S)\geq0.
\end{equation}

$S$ は任意なので命題が成立する。
\end{proof}

特に、

\begin{equation}
    B,I,K\in[0,1]
\end{equation}

ならば、

\begin{equation}
    0\leq BIK\leq1
\end{equation}

である。

%------------------------------------------------------------

\subsection{C.4 Solver 集合上の測度}

測度

\begin{equation}
    \mu_{\Solver}:
    \Sigma_{\Solver}\to[0,\infty]
\end{equation}

を導入する。

測度の可算加法性より、互いに素な
$A_n\in\Sigma_{\Solver}$ に対して

\begin{equation}
    \mu_{\Solver}
    \left(
        \bigcup_{n=1}^{\infty}A_n
    \right)
    =
    \sum_{n=1}^{\infty}
    \mu_{\Solver}(A_n)
\end{equation}

である。

ここでも、

\begin{equation}
    \Sigma_{\Solver}
\end{equation}

の存在と、

\begin{equation}
    \mu_{\Solver}
\end{equation}

の存在は別問題である。

%------------------------------------------------------------

\subsection{C.5 正規化}

\begin{equation}
    Z
    :=
    \int_{\Solver}
    g(S)\,d\mu_{\Solver}(S)
\end{equation}

と置く。

\begin{definition}[正規化可能性]
$g$ が正規化可能であるとは、

\begin{equation}
    0<Z<\infty
\end{equation}

が成立することをいう。
\end{definition}

このとき、

\begin{equation}
    \widehat g(S)
    :=
    \frac{g(S)}{Z}
\end{equation}

と定義する。

\begin{proposition}[正規化]
$g$ が非負可測関数であり、

\begin{equation}
    0<Z<\infty
\end{equation}

ならば、

\begin{equation}
    \int_{\Solver}
    \widehat g(S)\,d\mu_{\Solver}(S)
    =
    1
\end{equation}

である。
\end{proposition}

\begin{proof}
$0<Z<\infty$ より $1/Z$ は有限の正数である。

したがって積分の線形性から、

\begin{align}
    \int_{\Solver}
    \widehat g(S)\,d\mu_{\Solver}(S)
    &=
    \frac1Z
    \int_{\Solver}
    g(S)\,d\mu_{\Solver}(S)
    \\
    &=
    \frac ZZ
    \\
    &=1.
\end{align}

\end{proof}

%------------------------------------------------------------

\subsection{C.6 離散形式}

有限または可算な集合上で

\begin{equation}
    w:S\to[0,\infty)
\end{equation}

が与えられているならば、

\begin{equation}
    \widetilde\mu(i)
    =
    \sum_{b\in B(i)}w(b)
\end{equation}

と定義できる。

さらに

\begin{equation}
    Z
    :=
    \sum_{u\in S}w(u)
\end{equation}

が

\begin{equation}
    0<Z<\infty
\end{equation}

ならば、

\begin{equation}
    \mu(i)
    :=
    \frac{\widetilde\mu(i)}{Z}
\end{equation}

によって正規化できる。

実際、

\begin{equation}
    0
    \leq
    \widetilde\mu(i)
    \leq
    Z
\end{equation}

なので、

\begin{equation}
    0\leq\mu(i)\leq1
\end{equation}

である。

%------------------------------------------------------------

\subsection{C.7 連続形式}

可測空間

\begin{equation}
    (S,\Sigma_S,\mu_S)
\end{equation}

が与えられ、$i\in\Sigma_S$ ならば、

\begin{equation}
    \mu_S(i)
    =
    \int_S
    \mathbf1_i(u)\,d\mu_S(u)
\end{equation}

である。

これは指示関数の積分による通常の測度の表現である。

したがって、ECA の離散形式と連続形式を対応させる際には、
少なくとも

\begin{equation}
    i\in\Sigma_S
\end{equation}

という可測性が必要である。

%------------------------------------------------------------

\subsection{C.8 測度そのものと測度空間}

測度そのものを対象とする場合には、

\begin{equation}
    C^{(\mu)}=\mu
\end{equation}

を対象とする。

一方、測度空間を対象とする場合には、

\begin{equation}
    C^{(\mathrm{ms})}
    =
    (X,\Sigma_X,\mu_X)
\end{equation}

を対象とする。

一般には、

\begin{equation}
    \mu
\end{equation}

だけから

\begin{equation}
    (X,\Sigma_X)
\end{equation}

を一意に復元することはできない。

したがって、

\begin{equation}
    C^{(\mu)}
    \neq
    C^{(\mathrm{ms})}
\end{equation}

として扱う必要がある。

%################################################################
% (D)
%################################################################

\section{(D) 一般解の存在}

\subsection{D.1 一般解存在命題}

一般解の存在を

\begin{equation}
    \exists\,
    \Eval:\Solver\to X
\end{equation}

と定式化する。

これは単に各 Solver に個別の値が存在するという意味ではなく、
Solver 集合全体に対して同一の写像として評価構造を与えることを意味する。

%------------------------------------------------------------

\subsection{D.2 評価写像からの存在}

\begin{proposition}[評価写像による一般解の存在]
写像

\begin{equation}
    \Phi:\Solver\to X
\end{equation}

が存在し、これを ECA の要求する評価構造として採用できるならば、

\begin{equation}
    \Eval:=\Phi
\end{equation}

によって一般解が存在する。
\end{proposition}

\begin{proof}
$\Phi$ は仮定により

\begin{equation}
    \Phi:\Solver\to X
\end{equation}

という写像である。

したがって

\begin{equation}
    \Eval:=\Phi
\end{equation}

と定義すれば、

\begin{equation}
    \Eval:\Solver\to X
\end{equation}

が得られる。

これは一般解の存在命題そのものである。
\end{proof}

\begin{remark}
したがって、(D) の存在証明において本質的なのは、
評価写像 $\Phi$ が実際に構成できることを示すことである。

(A)--(C) の定義だけから $\Phi$ の存在が自動的に従うわけではない。
\end{remark}

%################################################################
% (E)
%################################################################

\section{(E) Solver 同型と一般解の構造的一意性}

\subsection{E.1 Solver 同型}

二つの Solver 表現

\begin{equation}
    \Solver_1,\qquad\Solver_2
\end{equation}

の間に Solver 同型

\begin{equation}
    T:
    \Solver_1
    \overset{\sim}{\longrightarrow}
    \Solver_2
\end{equation}

があるとする。

ここで $T$ は単なる集合論的全単射ではなく、
(A) で定義された Solver 構造を保存する写像である。

%------------------------------------------------------------

\subsection{E.2 評価空間間の構造保存写像}

一般に

\begin{equation}
    \Phi_1:\Solver_1\to X_1,
    \qquad
    \Phi_2:\Solver_2\to X_2
\end{equation}

とする。

評価空間が異なる場合には、

\begin{equation}
    U:X_1\to X_2
\end{equation}

を考える。

一般解が Solver 同型に対して構造保存的であることを、

\begin{equation}
    \boxed{
        U\circ\Phi_1
        =
        \Phi_2\circ T
    }
\end{equation}

によって表す。

%------------------------------------------------------------

\subsection{E.3 構造的一意性から導かれる最終評価の一致}

最終評価写像を

\begin{equation}
    \Psi_1:X_1\to\R,
    \qquad
    \Psi_2:X_2\to\R
\end{equation}

とする。

さらに、

\begin{equation}
    \Psi_2\circ U
    =
    \Psi_1
\end{equation}

が成立すると仮定する。

\begin{theorem}[Solver 構造に基づく一般解の構造的一意性]
上記の仮定のもとで、

\begin{equation}
    g_1:=\Psi_1\circ\Phi_1,
    \qquad
    g_2:=\Psi_2\circ\Phi_2
\end{equation}

と定義すれば、

\begin{equation}
    \boxed{
        g_2\circ T=g_1
    }
\end{equation}

が成立する。
\end{theorem}

\begin{proof}
合成写像を直接計算する。

\begin{align}
    g_2\circ T
    &=
    \Psi_2\circ\Phi_2\circ T
    \\
    &=
    \Psi_2\circ U\circ\Phi_1
    \\
    &=
    \Psi_1\circ\Phi_1
    \\
    &=
    g_1.
\end{align}

第2の等号には

\begin{equation}
    \Phi_2\circ T
    =
    U\circ\Phi_1
\end{equation}

を用い、第3の等号には

\begin{equation}
    \Psi_2\circ U
    =
    \Psi_1
\end{equation}

を用いた。

よって、

\begin{equation}
    g_2\circ T=g_1
\end{equation}

である。
\end{proof}

%------------------------------------------------------------

\subsection{E.4 共通評価空間の場合}

特に

\begin{equation}
    X_1=X_2=X
\end{equation}

かつ

\begin{equation}
    U=\operatorname{id}_X
\end{equation}

ならば、

\begin{equation}
    \Phi_2\circ T=\Phi_1
\end{equation}

となる。

したがって、

\begin{equation}
    \Phi_2(T(S))
    =
    \Phi_1(S)
    \qquad
    (\forall S\in\Solver_1).
\end{equation}

%------------------------------------------------------------

\subsection{E.5 証明可能性の限界}

ここで注意すべき点がある。

\begin{equation}
    U\circ\Phi_1
    =
    \Phi_2\circ T
\end{equation}

という可換性そのものは、
(A)--(D) の定義のみから自動的には導出されない。

これは Solver 同型 $T$ が保存する構造と、
評価写像 $\Phi_1,\Phi_2$ がどのようにその構造を反映するかについての
追加条件である。

したがって、現段階では、

\begin{equation}
    \boxed{
    \text{可換性}
    \Longrightarrow
    \text{構造的一意性}
    }
\end{equation}

は厳密に証明できるが、

\begin{equation}
    \boxed{
    \text{Solver 構造}
    \Longrightarrow
    \text{可換性}
    }
\end{equation}

については別途 ECA 固有の構造を用いた証明が必要である。

この点は、現在の定義体系における未証明部分として明示する。

%################################################################
% (F)
%################################################################

\section{(F) 離散 Solver の細分化極限}

\subsection{F.1 測度分割}

可測空間

\begin{equation}
    (I,\mathcal F,\mu)
\end{equation}

を考え、

\begin{equation}
    \mu(I)=1
\end{equation}

とする。

有限段階 $N$ の可測分割を

\begin{equation}
    \mathcal P_N
    =
    \{E_1^{(N)},\ldots,E_N^{(N)}\}
\end{equation}

とし、

\begin{equation}
    E_j^{(N)}\in\mathcal F,
\end{equation}

\begin{equation}
    E_p^{(N)}
    \cap
    E_q^{(N)}
    =
    \varnothing
    \qquad(p\neq q),
\end{equation}

\begin{equation}
    \bigcup_{j=1}^{N}E_j^{(N)}
    =
    I
\end{equation}

を仮定する。

測度要素を

\begin{equation}
    \Delta\mu_j^{(N)}
    :=
    \mu(E_j^{(N)})
\end{equation}

とする。

すると、

\begin{equation}
    \sum_{j=1}^{N}
    \Delta\mu_j^{(N)}
    =
    \mu(I)
    =
    1
\end{equation}

である。

%------------------------------------------------------------

\subsection{F.2 測度メッシュ}

\begin{equation}
    \left\|\mathcal P_N\right\|_\mu
    :=
    \max_{1\leq j\leq N}
    \Delta\mu_j^{(N)}
\end{equation}

と定義する。

細分化条件を

\begin{equation}
    \boxed{
        \left\|\mathcal P_N\right\|_\mu
        \longrightarrow0
        \qquad(N\to\infty)
    }
\end{equation}

とする。

これは単なる $N\to\infty$ ではなく、
各分割要素の測度が一様に小さくなることを意味する。

%------------------------------------------------------------

\subsection{F.3 離散評価}

代表点

\begin{equation}
    x_j^{(N)}\in E_j^{(N)}
\end{equation}

を選ぶ。

\begin{equation}
    f(x)
    :=
    B(x)I(x)K(x)
\end{equation}

とおけば、

\begin{equation}
    T_N
    :=
    \sum_{j=1}^{N}
    f(x_j^{(N)})
    \Delta\mu_j^{(N)}
\end{equation}

を定義できる。

一様正規化の場合には

\begin{equation}
    \Delta\mu_j^{(N)}
    =
    \frac1N
\end{equation}

であるから、

\begin{equation}
    \boxed{
        T_N
        =
        \frac1N
        \sum_{j=1}^{N}
        B(x_j^{(N)})
        I(x_j^{(N)})
        K(x_j^{(N)})
    }
\end{equation}

となる。

これは ECA の既存離散式をそのまま回収する。

%------------------------------------------------------------

\subsection{F.4 標準的な極限定理}

\begin{theorem}[測度付き離散評価の極限]
$f$ が適切な測度空間上で積分可能であり、
分割列 $\{\mathcal P_N\}$ が対応する積分の近似列として収束するならば、

\begin{equation}
    \lim_{N\to\infty}
    T_N
    =
    \int_I f(x)\,d\mu(x)
\end{equation}

である。
\end{theorem}

\begin{proof}
仮定より、

\begin{equation}
    T_N
    =
    \sum_{j=1}^{N}
    f(x_j^{(N)})
    \Delta\mu_j^{(N)}
\end{equation}

は $\int_I f\,d\mu$ の積分近似列である。

さらに、その近似列が積分に収束することを仮定しているため、

\begin{equation}
    \lim_{N\to\infty}T_N
    =
    \int_I f(x)\,d\mu(x)
\end{equation}

が成立する。
\end{proof}

$ f=B I K $ を代入すれば、

\begin{equation}
    \boxed{
        \lim_{N\to\infty}
        \sum_{j=1}^{N}
        B(x_j^{(N)})
        I(x_j^{(N)})
        K(x_j^{(N)})
        \Delta\mu_j^{(N)}
        =
        \int_I
        B(x)I(x)K(x)
        \,d\mu(x)
    }
\end{equation}

を得る。

%------------------------------------------------------------

\subsection{F.5 ECA の Solver 集合についての論考}

ここまでの定理が示しているのは、

\begin{equation}
    \text{適切な細分化}
    +
    \text{積分近似の収束条件}
    \Longrightarrow
    \text{連続形式への収束}
\end{equation}

である。

しかし、

\begin{equation}
    \text{ECA の任意の Solver 集合}
    \Longrightarrow
    \left\|\mathcal P_N\right\|_\mu\to0
\end{equation}

を現段階で証明したわけではない。

一様正規化の場合には、

\begin{equation}
    \left\|\mathcal P_N\right\|_\mu
    =
    \frac1N
    \longrightarrow0
\end{equation}

なので、測度メッシュの消失そのものは直ちに成立する。

一方、一般の測度では、

\begin{equation}
    \left\|\mathcal P_N\right\|_\mu
    =
    \max_j\mu(E_j^{(N)})
\end{equation}

であるから、

\begin{equation}
    N\to\infty
\end{equation}

だけでは十分ではない。

必要なのは、

\begin{equation}
    \max_j\mu(E_j^{(N)})
    \longrightarrow0
\end{equation}

である。

さらに、$BIK$ の可測性および可積分性も必要になる。

したがって、

\begin{equation}
    BIK\in L^1(\mu)
\end{equation}

等の条件を含め、
ECA の Solver 集合全体についてこの極限が必ず成立することは、
現段階では証明命題ではなく論考として位置付ける。

%################################################################
% (G)
%################################################################

\section{(G) Solver 集合における一般解の連続的構造}

\subsection{G.1 一般解の定義}

(A)--(F) の構造を前提として、

\begin{equation}
    \boxed{
        \Eval:\Solver\to X
    }
\end{equation}

を一般解とする。

各 $S\in\Solver$ に対して、

\begin{equation}
    \Eval(S)\in X
\end{equation}

である。

一般解は個々の Solver に対して個別の数値を列挙するものではなく、
Solver 集合全体に共通する評価構造を $X$ 上に表現する写像である。

%------------------------------------------------------------

\subsection{G.2 最終評価}

\begin{equation}
    \EvalMap:X\to\R
\end{equation}

とし、

\begin{equation}
    g
    =
    \EvalMap\circ\Eval
\end{equation}

と定義する。

したがって、

\begin{equation}
    \boxed{
        \Solver
        \xrightarrow{\Eval}
        X
        \xrightarrow{\EvalMap}
        \R
    }
\end{equation}

という二段階の構造が得られる。

%------------------------------------------------------------

\subsection{G.3 評価写像との関係}

評価写像

\begin{equation}
    \Phi:\Solver\to X
\end{equation}

と一般解

\begin{equation}
    \Eval:\Solver\to X
\end{equation}

は型が同じであっても、その意味を同一視する必要はない。

もし

\begin{equation}
    \Eval(S)=\Phi(S)
    \qquad
    (\forall S\in\Solver)
\end{equation}

が成立するならば、

\begin{equation}
    \Eval=\Phi
\end{equation}

として同一視できる。

しかしこれは一般解の型から自動的に従うものではない。

%------------------------------------------------------------

\subsection{G.4 (E) との接続}

Solver 同型

\begin{equation}
    T:\Solver_1\overset{\sim}{\longrightarrow}\Solver_2
\end{equation}

に対して、

\begin{equation}
    \Eval_2\circ T
    =
    \Eval_1
\end{equation}

が成立するならば、
(E) の構造的一意性と整合する。

ただし、この可換性そのものは
(E) と同様に追加の構造保存条件である。

%------------------------------------------------------------

\subsection{G.5 (F) との接続}

(F) の離散評価

\begin{equation}
    T_N
    =
    \sum_{j=1}^{N}
    f(x_j^{(N)})
    \Delta\mu_j^{(N)}
\end{equation}

が、

\begin{equation}
    \int_I f(x)\,d\mu(x)
\end{equation}

へ収束するならば、
一般解による評価構造は離散表現から連続表現へ接続される。

したがって、

\begin{equation}
    \boxed{
    \text{Solver}
    \longrightarrow
    \Eval
    \longrightarrow
    \text{離散評価}
    \longrightarrow
    \text{連続評価}
    }
\end{equation}

という構造が得られる。

%################################################################
% (H)
%################################################################

\section{(H) 中間 Solver の厳密な定義と構造}

\subsection{H.1 Solver--ECA 適合条件}

Solver

\begin{equation}
    s=(\Aset_s,\Rset_s)
\end{equation}

に対して ECA の公理選択空間

\begin{equation}
    I\subseteq\mathcal P(S)
\end{equation}

を考える。

\begin{definition}[Solver--ECA 適合条件]
Solver $s$ が ECA に適合するとは、

\begin{equation}
    \Aset_s\in I
\end{equation}

が成立することをいう。
\end{definition}

したがって、

\begin{equation}
    \Solver_I
    :=
    \left\{
        s\in\Solver
        \mid
        \Aset_s\in I
    \right\}
\end{equation}

と定義する。

%------------------------------------------------------------

\subsection{H.2 公理選択復元写像}

\begin{definition}[公理選択復元写像]
\begin{equation}
    \mathfrak A:
    \Solver_I\to I
\end{equation}

を

\begin{equation}
    \boxed{
        \mathfrak A(s)
        :=
        \Aset_s
    }
\end{equation}

によって定義する。

以下、

\begin{equation}
    i_s
    :=
    \mathfrak A(s)
    =
    \Aset_s
\end{equation}

と記す。
\end{definition}

\begin{proposition}[公理選択復元写像の well-defined 性]
$\mathfrak A$ は well-defined である。
\end{proposition}

\begin{proof}
$s\in\Solver_I$ を任意に取る。

$\Solver_I$ の定義より、

\begin{equation}
    \Aset_s\in I
\end{equation}

である。

また Solver は

\begin{equation}
    s=(\Aset_s,\Rset_s)
\end{equation}

として定義されているので、
同一の $s$ に対する $\Aset_s$ は一意に定まる。

したがって、

\begin{equation}
    \mathfrak A(s)=\Aset_s
\end{equation}

は $I$ の唯一の元として定まり、

\begin{equation}
    \mathfrak A:\Solver_I\to I
\end{equation}

は well-defined である。
\end{proof}

%------------------------------------------------------------

\subsection{H.3 理論写像との関係}

Solver の理論を

\begin{equation}
    \mathcal T(s)
    =
    \mathcal T(\Aset_s,\Rset_s)
\end{equation}

とする。

$ i_s=\Aset_s $ だから、

\begin{equation}
    \boxed{
        \mathcal T(s)
        =
        \mathcal T(i_s,\Rset_s)
    }
\end{equation}

である。

ここで重要なのは、

\begin{equation}
    \mathcal T(s)=i_s
\end{equation}

とはならないことである。

$\mathcal T(s)$ は理論であり、
$i_s$ はその Solver の公理選択である。

%------------------------------------------------------------

\subsection{H.4 離散評価構造}

ECA における

\begin{equation}
    B(i)\subseteq i\subseteq S
\end{equation}

を用いる。

\begin{definition}[Solver の離散評価構造]
$s\in\Solver_I$ に対して、

\begin{equation}
    \boxed{
        D_s
        :=
        B(i_s)
    }
\end{equation}

と定義する。
\end{definition}

したがって、

\begin{equation}
    D_s\subseteq i_s\subseteq S
\end{equation}

である。

重み関数

\begin{equation}
    w_s:S\to[0,1]
\end{equation}

が

\begin{equation}
    \sum_{u\in S}w_s(u)=1
\end{equation}

を満たす場合、

\begin{equation}
    \boxed{
        E_{\mathrm d}(s)
        :=
        \sum_{b\in D_s}w_s(b)
    }
\end{equation}

と定義する。

%------------------------------------------------------------

\subsection{H.5 離散評価の有界性}

\begin{proposition}
\begin{equation}
    0\leq E_{\mathrm d}(s)\leq1
\end{equation}

である。
\end{proposition}

\begin{proof}
各 $b\in D_s$ に対して

\begin{equation}
    w_s(b)\geq0
\end{equation}

なので、

\begin{equation}
    E_{\mathrm d}(s)\geq0.
\end{equation}

また $D_s\subseteq S$ より、

\begin{equation}
    \sum_{b\in D_s}w_s(b)
    \leq
    \sum_{u\in S}w_s(u)
    =
    1.
\end{equation}

したがって、

\begin{equation}
    0\leq E_{\mathrm d}(s)\leq1.
\end{equation}

\end{proof}

%------------------------------------------------------------

\subsection{H.6 連続評価構造：測度そのもの}

Solver $s$ に対応する測度を $\mu_s$ とする。

\begin{definition}[連続評価測度]
\begin{equation}
    \boxed{
        C_s^{(\mu)}
        :=
        \mu_s
    }
\end{equation}

と定義する。
\end{definition}

$i_s$ が可測集合であるなら、

\begin{equation}
    E_{\mathrm c}(s)
    :=
    \mu_s(i_s)
    =
    \int_S
    \mathbf1_{i_s}(u)
    \,d\mu_s(u)
\end{equation}

である。

%------------------------------------------------------------

\subsection{H.7 連続評価構造：測度空間}

測度空間全体を対象とする場合には、

\begin{definition}[連続評価測度空間]
\begin{equation}
    \boxed{
        C_s^{(\mathrm{ms})}
        :=
        (I_s,\Sigma_{I_s},\mu_s)
    }
\end{equation}

とする。
\end{definition}

ここで

\begin{equation}
    \Sigma_{I_s}
\end{equation}

は $I_s$ 上の $\sigma$-代数であり、

\begin{equation}
    \mu_s:
    \Sigma_{I_s}\to[0,1]
\end{equation}

は測度である。

したがって、

\begin{equation}
    C_s^{(\mu)}
\end{equation}

と

\begin{equation}
    C_s^{(\mathrm{ms})}
\end{equation}

は異なるレベルの対象である。

%------------------------------------------------------------

\subsection{H.8 測度空間から測度への射影}

\begin{equation}
    \pi_\mu:
    C_s^{(\mathrm{ms})}
    \to
    C_s^{(\mu)}
\end{equation}

を

\begin{equation}
    \boxed{
        \pi_\mu
        (I_s,\Sigma_{I_s},\mu_s)
        =
        \mu_s
    }
\end{equation}

によって定義する。

これは測度空間から測度成分だけを取り出す自然な忘却写像である。

逆に、一般には $\mu_s$ だけから

\begin{equation}
    (I_s,\Sigma_{I_s})
\end{equation}

を一意に復元することはできない。

%------------------------------------------------------------

\subsection{H.9 離散評価と連続評価の共通起点}

同じ公理選択 $i_s$ を用いることで、

\begin{equation}
    i_s
    \longrightarrow
    D_s=B(i_s)
    \longrightarrow
    E_{\mathrm d}(s)
\end{equation}

および

\begin{equation}
    i_s
    \longrightarrow
    \mu_s
    \longrightarrow
    E_{\mathrm c}(s)
\end{equation}

という二つの評価経路が得られる。

したがって、

\begin{equation}
    \boxed{
    \begin{aligned}
        i_s
        &\longrightarrow
        D_s
        \longrightarrow
        E_{\mathrm d}(s),
        \\[2mm]
        i_s
        &\longrightarrow
        C_s^{(\mu)}
        \longrightarrow
        E_{\mathrm c}(s)
    \end{aligned}
    }
\end{equation}

である。

重要なのは、離散評価と連続評価が
別々の公理選択から発生するのではなく、
同じ $i_s$ を起点として構成されることである。

%------------------------------------------------------------

\subsection{H.10 中間 Solver の厳密な定義}

Solver $s\in\Solver_I$ に対して、

\begin{equation}
    \Eval(s)
    =
    \Phi_s
    \left(
        E_{\mathrm d}(s),
        E_{\mathrm c}(s)
    \right)
\end{equation}

と表現できる場合を考える。

\begin{definition}[中間 Solver]
$s$ が中間 Solver であるとは、

\begin{enumerate}[label=\textup{(\arabic*)}]

\item

\begin{equation}
    i_s=\mathfrak A(s)=\Aset_s\in I
\end{equation}

である。

\item

\begin{equation}
    D_s=B(i_s)
\end{equation}

および

\begin{equation}
    E_{\mathrm d}(s)
    =
    \sum_{b\in D_s}w_s(b)
\end{equation}

が定義される。

\item

\begin{equation}
    C_s^{(\mu)}=\mu_s
\end{equation}

または、必要に応じて

\begin{equation}
    C_s^{(\mathrm{ms})}
    =
    (I_s,\Sigma_{I_s},\mu_s)
\end{equation}

が定義される。

\item

\begin{equation}
    \Eval(s)
    =
    \Phi_s
    \left(
        E_{\mathrm d}(s),
        E_{\mathrm c}(s)
    \right)
\end{equation}

であり、$\Phi_s$ が第1変数および第2変数の双方に
非自明に依存する。

\end{enumerate}
\end{definition}

最後の条件は、

\begin{equation}
    \exists x_1\neq x_2,\ y:
    \quad
    \Phi_s(x_1,y)
    \neq
    \Phi_s(x_2,y)
\end{equation}

および

\begin{equation}
    \exists x,\ y_1\neq y_2:
    \quad
    \Phi_s(x,y_1)
    \neq
    \Phi_s(x,y_2)
\end{equation}

によって厳密に表現できる。

この定義では、中間 Solver を

\begin{equation}
    \text{離散構造}\cup\text{連続構造}
\end{equation}

として定義していない。

必要なのは、

\begin{equation}
    \boxed{
    \text{同一 Solver 内に双方の評価構造が存在し、
    その双方が一般解の決定に非自明に寄与する}
    }
\end{equation}

ことである。

%------------------------------------------------------------

\subsection{H.11 中間 Solver 構造定理}

\begin{theorem}[中間 Solver 構造定理]
\label{thm:middle-solver}

$s\in\Solver_I$ が

\begin{equation}
    \Eval(s)
    =
    \Phi_s(E_{\mathrm d}(s),E_{\mathrm c}(s))
\end{equation}

と表現できるとする。

このとき、$s$ が中間 Solver であることは、

\begin{equation}
    \exists x_1\neq x_2,\ y:
    \quad
    \Phi_s(x_1,y)
    \neq
    \Phi_s(x_2,y)
\end{equation}

かつ

\begin{equation}
    \exists x,\ y_1\neq y_2:
    \quad
    \Phi_s(x,y_1)
    \neq
    \Phi_s(x,y_2)
\end{equation}

が成立することと同値である。
\end{theorem}

\begin{proof}
中間 Solver の定義より、
$\Phi_s$ は第1変数および第2変数の双方に
非自明に依存しなければならない。

第1変数について、もし

\begin{equation}
    \Phi_s(x_1,y)
    =
    \Phi_s(x_2,y)
    \qquad
    (\forall x_1,x_2,y)
\end{equation}

ならば、$\Phi_s$ は第1変数に依存しない。

実際、任意の $y$ を固定すると、
$\Phi_s(x,y)$ は $x$ に依存せず一定であるから、
ある写像

\begin{equation}
    \Psi_s:[0,1]\to X
\end{equation}

を用いて

\begin{equation}
    \Phi_s(x,y)=\Psi_s(y)
\end{equation}

と書ける。

したがって、

\begin{equation}
    \Eval(s)
    =
    \Psi_s(E_{\mathrm c}(s))
\end{equation}

となり、一般解は離散評価に依存しない。

これは中間 Solver の定義に反する。

よって、

\begin{equation}
    \exists x_1\neq x_2,\ y:
    \Phi_s(x_1,y)\neq\Phi_s(x_2,y)
\end{equation}

が必要である。

第2変数についても同様に、
もし

\begin{equation}
    \Phi_s(x,y_1)
    =
    \Phi_s(x,y_2)
    \qquad
    (\forall x,y_1,y_2)
\end{equation}

ならば、$\Phi_s$ は第2変数に依存しない。

したがって、

\begin{equation}
    \exists x,\ y_1\neq y_2:
    \Phi_s(x,y_1)\neq\Phi_s(x,y_2)
\end{equation}

が必要である。

逆に、両方の非自明な依存性が成立すれば、
中間 Solver の第4条件を満たす。

よって同値性が成立する。
\end{proof}

%------------------------------------------------------------

\subsection{H.12 $\mathcal{T}(s)$ から評価構造への接続}

Solver

\begin{equation}
    s=(\Aset_s,\Rset_s)
\end{equation}

に対して、

\begin{equation}
    i_s=\Aset_s
\end{equation}

であるから、

\begin{equation}
    \mathcal T(s)
    =
    \mathcal T(i_s,\Rset_s)
\end{equation}

である。

さらに、

\begin{equation}
    D_s=B(i_s)
\end{equation}

なので、

\begin{equation}
    i_s
    \longrightarrow
    D_s
    \longrightarrow
    E_{\mathrm d}(s)
\end{equation}

という経路を得る。

同様に、

\begin{equation}
    i_s
    \longrightarrow
    \mu_s
    \longrightarrow
    E_{\mathrm c}(s)
\end{equation}

という経路を得る。

したがって、

\begin{equation}
    \boxed{
        s
        \longrightarrow
        i_s
        \longrightarrow
        \left(
            E_{\mathrm d}(s),
            E_{\mathrm c}(s)
        \right)
        \longrightarrow
        \Eval(s)
    }
\end{equation}

という構造が成立する。

%------------------------------------------------------------

\subsection{H.13 最終的な構造}

以上の結果をまとめると、

\begin{equation}
    s=(\Aset_s,\Rset_s)
\end{equation}

から、

\begin{equation}
    i_s
    =
    \mathfrak A(s)
    =
    \Aset_s
\end{equation}

を得る。

さらに、

\begin{equation}
    D_s=B(i_s)
\end{equation}

および

\begin{equation}
    C_s^{(\mu)}=\mu_s
\end{equation}

または

\begin{equation}
    C_s^{(\mathrm{ms})}
    =
    (I_s,\Sigma_{I_s},\mu_s)
\end{equation}

を得る。

そして、

\begin{equation}
    E_{\mathrm d}(s)
    =
    \sum_{b\in D_s}w_s(b)
\end{equation}

および

\begin{equation}
    E_{\mathrm c}(s)
    =
    \int_S
    \mathbf1_{i_s}(u)
    \,d\mu_s(u)
\end{equation}

を構成する。

したがって、

\begin{equation}
    \boxed{
        \text{Solver}
        \longrightarrow
        \text{公理選択}
        \longrightarrow
        \text{離散・連続評価構造}
        \longrightarrow
        \text{一般解}
    }
\end{equation}

という構造が得られる。

特に中間 Solver は、

\begin{equation}
    \boxed{
    \begin{minipage}{0.85\linewidth}
    同一の Solver に対して離散的評価構造と連続的評価構造が存在し、
    その双方が一般解の決定に非自明に寄与する Solver
    \end{minipage}
    }
\end{equation}

として定義される。

%################################################################
% 総括
%################################################################

\section{総括：証明可能事項と追加条件}

本稿において、改訂版 v2 の定義体系から直接または
標準的な数学的事実によって導出できる主要な結果は次の通りである。

\begin{enumerate}

\item
部分 Solver の包含関係から、

\begin{equation}
    S'\SubSolver S
    \land
    |S'|>\aleph_0
    \Longrightarrow
    |S|>\aleph_0
\end{equation}

が成立する。

\item
可測な一般解と可測な最終評価写像から、

\begin{equation}
    g=\EvalMap\circ\Eval
\end{equation}

の可測性が成立する。

\item
\[
    0<
    \int_{\Solver}g\,d\mu_{\Solver}
    <\infty
\]
ならば、

\begin{equation}
    \widehat g
    =
    \frac{g}{Z}
\end{equation}

は正規化される。

\item
評価写像が実際に構成されれば、

\begin{equation}
    \Eval:\Solver\to X
\end{equation}

として一般解が存在する。

\item
Solver 同型および評価空間間の構造保存写像が

\begin{equation}
    U\circ\Phi_1
    =
    \Phi_2\circ T
\end{equation}

を満たし、さらに最終評価を保存するならば、

\begin{equation}
    g_2\circ T=g_1
\end{equation}

が成立する。

\item
適切な細分化条件と積分近似の収束条件のもとで、

\begin{equation}
    T_N
    \longrightarrow
    \int_I B(x)I(x)K(x)\,d\mu(x)
\end{equation}

が成立する。

\item
ECA に適合する Solver について、

\begin{equation}
    i_s
    =
    \mathfrak A(s)
    =
    \Aset_s
\end{equation}

は well-defined である。

\item
その結果、

\begin{equation}
    D_s=B(i_s)
\end{equation}

および連続評価構造 $C_s$ を
同一の公理選択 $i_s$ から構成できる。

\item
一般解が離散評価と連続評価の双方に非自明に依存することと、
中間 Solver の定義は同値である。

\end{enumerate}

一方、現段階で追加の ECA 固有の証明を必要とするものは、

\begin{equation}
    \boxed{
    \text{Solver 構造}
    \Longrightarrow
    U\circ\Phi_1=\Phi_2\circ T
    }
\end{equation}

という可換性そのもの、および

\begin{equation}
    \boxed{
    \text{任意の ECA Solver 集合}
    \Longrightarrow
    \left\|\mathcal P_N\right\|_\mu\to0
    }
\end{equation}

という細分化極限の一般成立である。

したがって、現時点の ECA の論理構造は、

\begin{equation}
\boxed{
\begin{array}{c}
\text{(A) Solver 構造}\\
\Downarrow\\
\text{(B) 一般評価空間}\\
\Downarrow\\
\text{(C) 測度論的評価}\\
\Downarrow\\
\text{(D) 一般解の存在}\\
\Downarrow\\
\text{(E) Solver 構造的一意性}\\
\Downarrow\\
\text{(F) 離散--連続極限}\\
\Downarrow\\
\text{(G) 一般解の連続的構造}\\
\Downarrow\\
\text{(H) 中間 Solver}
\end{array}
}
\end{equation}

として整理できる。

特に (H) では、

\begin{equation}
    \boxed{
        s
        \longrightarrow
        i_s
        \longrightarrow
        (D_s,C_s)
        \longrightarrow
        \left(
            E_{\mathrm d}(s),
            E_{\mathrm c}(s)
        \right)
        \longrightarrow
        \Eval(s)
    }
\end{equation}

という具体的な導出経路が確立されており、
(A)--(G) で定義された構造が、
(H) において Solver の内部構造へ接続される。

\end{document}
